\chapter{Eine Zurechtstellung und Erklärung}
\quellseite{5}

Leider sehen wir uns genöthigt, noch einmal auf P. Gößwein und seine Anklagen gegen unsere Synode und die hiesige Facultät zurückzukommen. P. G. fährt nämlich fort, solche Anklagen in die Welt hinaus zu schreiben. Er hat angefangen, ein monatliches Blättchen herauszugeben, welches gerade diesem Zweck dienen soll, wie er sich ausdrückt, „Zeugniß“ abzulegen gegen „Neumissouri“. Auch in diesem Blatt ergeht sich P. G. in Schmähungen, in Verdrehungen und Entstellungen von Worten und Thatsachen, ja, in offenbaren Erdichtungen und Einbildungen. Es ist schier unmöglich, ist auch nicht nöthig, alles, was P. G. schreibt, zurechtzustellen und zu widerlegen. Wir heben hier nur einzelne Punkte heraus, die P. G.’s Art und Weise ganz besonders klar stellen.

1. P. G. legt einzelnen Gliedern der Facultät fort und fort Worte in den Mund, die sie nie gesprochen, schiebt ihnen fort und fort Lehren unter, die sie nie gelehrt haben. So heißt es z. B. in seinem Blatt: „Nun heißen Dr. Gräbner und seine Freunde die Kirche vor Christi zweitem Advent am jüngsten Tage auf einen Mann warten, welcher ein biblisches Buch nach Text und Auslegung durch unmittelbare göttliche Erleuchtung erst bekräftigen muß, ehe sie etwas Gewisses daraus entnehmen darf.“ (S. 2.) Ferner: „Dr. Gräbner aber und Genossen sagen der Kirche, der Heilige Geist könne ihr auch ein Buch gegeben haben, das ihr völlig unnütz bleiben soll, solange er es durch seine himmlische Mission deutet und beglaubigt.“ (S. 2.) „So spricht Dr. Gräbner: ‚Mir war je und je nichts so zuwider als Walthers und seiner Freunde Festigkeit und Zuversichtlichkeit wider ihre Gegner. Text und Auslegung müssen ungewiß bleiben. Es kann nichts Gewisses geben.‘“ (S. 3.) „Das können sie aber nicht, wenn sie auf ihre Autorität pochen, wie Präses Pieper, der auf meine Frage: ‚Habe ich Sie persönlich beleidigt?‘ antwortete: ‚Das kann ich nicht sagen, aber bedenken Sie, Sie haben mir gegenüber auch kein Blatt vor den Mund genommen.‘“ (S. 7.) Alle diese Ausdrücke und noch viele andere, die P. G. in seinen Publikationen Gliedern der Facultät in den Mund gelegt hat, sind entweder von diesen Männern nie gebraucht, oder sie sind so entstellt, aus dem Zusammenhang gerissen und in einen anderen Zusammenhang gebracht, daß die Worte nun ganz anders lauten, als sie eigentlich gemeint waren.

2. Allerdings versucht P. G. in seinem Blatte einen Beweis für seine Behauptungen zu erbringen. Er citirt zu dem Zweck wiederholt „Lehre und Wehre“, Jahrg. 47, S. 263. Auch hier zeigt sich wieder P. G.’s Art und Weise. Einmal sagen die angeführten Worte das gar nicht, was P. G. ihnen entnimmt, daß die Offenbarung Johannis nach Text und Auslegung ungewiß bleiben müsse, daß sie völlig unnütz sei; sie heißen die Kirche nicht
\quellseite{6}
auf einen neuen Propheten warten. Sodann aber sind diese Worte auch gar nicht Worte, die Dr. Gräbner oder seine Freunde und Genossen gesagt oder geschrieben haben, wie es nach P. G.’s Citat den Anschein gewinnen muß, sondern diese Worte finden sich im 9. Bericht unserer Allgemeinen Synode vom Jahre 1857 und sind zu dem Zweck in „Lehre und Wehre“ aufgenommen, um auch an ihnen mit zu zeigen und nachzuweisen, was unsere Väter je und je von dem 20. Capitel der Offenbarung und dessen Auslegung gelehrt und bezeugt haben.

3. Ferner schreibt P. G. in seinem Blatt: „So verführt der Irrstern von St. Louis die Schriften zur Menschenrechtschaft und macht es ihnen zur Pflicht, jede Klageschrift gegen die Fakultät ungeprüft zu lassen, sich vielmehr derselben anzunehmen, wie Juden und Logenbrüder den Ihrigen durchhelfen.“ Das ist offenbare Unwahrheit. Die Facultät hat es nie irgend einem Gliede der Synode zur Pflicht gemacht, jede oder auch nur irgend eine Klageschrift gegen die Facultät ungeprüft zu lassen. P. G. wird den Beweis für diese Anklage wohl schuldig bleiben. P. G.’s Klageschrift gegen die Facultät ist von der Synode in genau derselben Weise geprüft worden, wie die Synode jede andere Klageschrift gegen eins ihrer Glieder oder eine ihrer Behörden prüft.

4. Endlich schreibt P. G. in einem Artikel im „Kirchl. Sendbote“ Folgendes: „Der Allgemeine Präses wurde wiederholt wegen Tyrannei angeklagt und nach einer solchen Synodalverhandlung sprach sein eigener Bruder in Gegenwart des Prof. Herzer zu mir: ‚Wehe der armen Seele meines Bruders, wenn Alles wahr ist, was da N. N. von ihm ausgesagt hat! Ich woll’s nicht hoffen, aber auch nicht bestreiten. Die Präseswürde ist zu gefährlich für einen theologischen Professor. Er wird nur allzu leicht versucht, Synodale für Untergebene anzusehen. Ich habe meinen Bruder schon dringend gebeten, das Präsidium niederzulegen; denn ich getraute mich an seiner Stelle nicht selig zu werden. Wie oft habe ich Gott schon gedankt, daß er mich nicht in gleiche Lage wie ihn kommen ließ! Ach, wie schwer kann die Versuchung werden unter Schmeichlern, die Alles gut heißen.‘“

Dieser ganze Passus enthält nicht etwa Entstellungen und Verdrehungen der Wahrheit, sondern ganz einfach eitel Erdichtungen und Einbildungen. Dr. Pieper ist nie auf der Synode der Tyrannei angeklagt worden, wie die ganze Synode bezeugen kann. Es hat also nie eine „solche Synodalverhandlung“ stattgefunden, nach der jemand zu P. G. etwas Derartiges hätte sagen können. Die Worte, die P. G. dem „eigenen Bruder“ Dr. Piepers in den Mund legt, sind nie gesprochen worden. Wir haben uns an die drei Brüder Dr. Piepers gewandt, die im Amte stehen, ob sie sich zu diesen Worten bekennen. Prof. Aug. Pieper und P. N. Pieper schreiben, daß sie, soweit sie sich erinnern können, mit P. G. seit Ablauf der gemeinsamen Studienzeit, also seit dem Jahre 1879, überhaupt nicht mehr gesprochen hätten. Prof. R. Pieper, der wahrscheinlich mit dem „eigenen Bruder“ gemeint ist, schreibt: „Wie ich Ihnen schon mündlich mittheilte, habe ich Herrn P. Gößwein nur einmal gesehen und gesprochen, und zwar auf meiner Rückkehr von der Versammlung der Synodalconferenz in Evansville, Ind., vor etwa sechs bis acht Jahren. Ich war nur kurze Zeit bei ihm, und meines Wissens ist von uns über die Stellung des Allgemeinen Präses unserer Synode überhaupt nicht geredet worden. Damals dachte auch keiner daran, daß mein Bruder jemals zum Allgemeinen Präses gewählt werden würde. Da er aber damals dieses Amt noch nicht bekleidete, so geht schon daraus hervor, daß ich die mir von Herrn P. Gößwein in den Mund gelegten Worte unmöglich geredet haben kann. Dieselben sind jedenfalls ein Erzeugniß seiner krankhaft erregten Phantasie, was bei dem geistig-körperlichen Zustande Herrn P. Gößweins leicht erklärlich ist. Seit der Bruder Allgemeiner Präses ist, habe ich Herrn P. Gößwein weder gesehen noch auch nur eine Zeile mit ihm gewechselt, soweit meine Erinnerung reicht.“

Was soll man nun zu solchen Dingen sagen, wenn man auch das mit hinzunimmt, was der „Lutheraner“ in No. 25 des letzten Jahrgangs, S. 390 f., dargethan hat? Denn jene Unwahrheiten und Entstellungen, die wir dort aufgedeckt haben, gehören auch auf P. G.’s Rechnung. Das „Kirchenblatt“ hat, wie wir nachträglich erfahren, seine Angaben aus dem eben angeführten Artikel P. G.’s im „Kirchl. Sendboten“ genommen. Wir fragen, was soll man zu diesem Allen sagen? Es ist da nur ein Entweder — oder möglich. Entweder ist P. G. ein Mann, der wider besseres Wissen und Gewissen die Unwahrheit sagt, ein Mann, der mit groben, schändlichen Lügen umgeht und sich nicht scheut, rein selbsterfundene und -erdachte Dinge boshafter Weise in die Welt hinaus zu schreiben. Oder aber P. G. ist in einem solchen krankhaft erregten Zustande, daß er nicht mehr weiß, was er sagt, daß seine Phantasie ihm in krankhafter Weise Begebenheiten als geschehen und Worte als gesprochen vorspiegelt, die nie geschehen und gesprochen sind. Die Facultät hat stets das Letztere angenommen und nimmt es heute noch an. Zu dieser Annahme bestimmt sie die Thatsache, daß P. G. schon seit langen Jahren, schon seit seiner Studienzeit, an einer schweren Krankheit, an der Epilepsie, leidet. Wir halten P. G. nicht im vollen Maße verantwortlich für das, was er sagt und schreibt.

Hiernach ist auch zu beurtheilen, was P. G. in seinem Blatte erzählt von dem Zeugenberuf, der ihm angeblich von dem seligen Dr. Walther aufgetragen sei. Soweit uns bekannt ist, weiß niemand in der ganzen Synode außer P. G. etwas von einer solchen Veranstaltung Dr. Walthers, wie sie P. G. beschreibt, bestimmten Personen die Ueberwachung einzelner Lehren zu übertragen. Auf diese Weise erklärt es sich auch, daß P. G. sich fortwährend für einen Verfolgten hält und erklärt, während in der That niemand ihn verfolgt, sondern, wenn von Verfolgen überhaupt geredet werden soll, dasselbe ganz auf seiner Seite ist.

Wir haben Grund zu der Vermuthung, daß P. G. fortfahren wird, in seiner Weise Synode und Facultät anzugreifen und zu verleumden und also seinen vermeintlichen Zeugenberuf auszurichten. Das können wir natürlich nicht hindern. Wir sind aber der Meinung, daß wir im „Lutheraner“ nicht weiter solche Auslassungen berücksichtigen sollten, die entweder grobe Lügen oder Einbildungen einer krankhaft erregten Phantasie sind.

G. M.
